% !PRINTED 648
\weberpage{648}
%<<CONT
tion into linear factors, then by 2. one of these linear factors must be divisible by $\pi$. Two different linear functions $z - c$ and $z - c'$ cannot be divisible by the same $\pi$, since otherwise the constant $c - c'$, which is different from zero, would be divisible by $\pi$.

6. \gesperrt{Every integral function $\omega$ in $\Omega$ is congruent to a constant $b$ with respect to the modulus $\pi$, i.\,e. the constant $b$ can be so chosen that
$\omega - b$ becomes divisible by $\pi$.}
% [sic] German source prints "das" for "daß".

For by \S\,175 the function $\omega$ satisfies an equation:
\[
\omega^n + a_1 \omega^{n-1} + \cdots + a_{n-1}\omega + a_n = 0,
\]
in which the $a_1, \ldots, a_n$ are integral rational functions of $z$. If $a_1^0, \ldots, a_n^0$ are the remainders of these functions on division by $z - c$, and hence constants, it follows that:
\[
\omega^n + a_1^0 \omega^{n-1} + \cdots + a_{n-1}^0 \omega + a_n^0 \equiv 0\ (\mathrm{mod}\ \pi),
\]
and if the function on the left side is decomposed into linear factors:
\weq{1}{(\omega - b)(\omega - b')(\omega - b'') \cdots \equiv 0\ (\mathrm{mod}\ \pi).}
At least one of these linear factors must therefore be divisible by $\pi$, which was to be proved.

From the one-valued decomposability of the functionals into prime factors that has now been established there follow far-reaching consequences, of which the most important shall be adduced here.

If in a functional $\varphi$ of the field $\Omega$ a certain number of the auxiliary variables, $x, y, \ldots$, occur only in the numerator, then
\[
\varphi = \varphi(x, y, \ldots)
\]
shall be called a \gesperrt{holomorphic function} of $x, y, \ldots$. Of these the theorem holds:

7. \gesperrt{A holomorphic function is an integral functional in $\overline{\Omega}$ only when the coefficients of the ordered function $\varphi$ are integral functionals.}

It is evidently sufficient to establish the theorem for holomorphic functions of one variable
\weq{2}{\varphi = x^m \varphi_0 + x^{m-1}\varphi_1 + \cdots + x \varphi_{m-1} + \varphi_m}
since from this it can be proved in general by complete induction.

This proof, however, follows from \gesperrt{Gauss's} theorem:
%CONT>>
